\begin{proof}[Proof of \cref{obj:thm:matched-risk-frontier}]
\begin{enumerate}
\item Fix \(n\ge2\) and \(0<\epsilon\le1\). The numerical benchmark satisfies
\[
0<r(n,\epsilon)\le1,
\]
because \(n^{-1/4}>0\) and the defining minimum is taken with \(1\).

Apply \cref{obj:lem:frontier-testing-priors}. Its probabilistic and combinatorial inputs are the bounded-differences exponential-moment bound from \citep[Theorem 6.2, displayed log-mgf conclusion in its proof, p.~167]{BoucheronLugosiMassart2013} and the labeled-tree count from \citep[Section 5.6, Theorem 5.40]{KellerTrotterAppliedCombinatorics}. Denote the resulting nonempty finite families and their uniformly averaged dataset laws by
\[
\begin{gathered}
N_{\mathrm{prior},0},N_{\mathrm{prior},1}\in\mathbb N_{\ge1},
\qquad
\mathcal F_\eta
=
(P_{\eta,\iota})_{\iota=1}^{N_{\mathrm{prior},\eta}},
\qquad \eta\in\{0,1\},\\
Q^{(\eta)}
=
\frac{1}{N_{\mathrm{prior},\eta}}
\sum_{\iota=1}^{N_{\mathrm{prior},\eta}}
(P_{\eta,\iota}^{\mathrm{obs}})^{\otimes n}.
\end{gathered}
\]
The supplied target separation \(\Delta\) obeys
\[
\Delta\ge2^{-14}r(n,\epsilon)>0,
\qquad
P_{\eta,\iota}\in\mathcal P,
\qquad
\theta(P_{\eta,\iota})=\eta\Delta.
\]
Moreover, for every \(M\in\mathfrak M_{n,\epsilon}(\mathbb R)\),
\[
d_{\mathrm{TV}}\!\left(MQ^{(0)},MQ^{(1)}\right)\le\frac14.
\]
% lean: CausalSmith.Stat.PrivateCateRoughdesign.frontier_testing_priors
% lean: CausalSmith.Stat.PrivateCateRoughdesign.rate_pos

\item Fix any \(M\in\mathfrak M_{n,\epsilon}(\mathbb R)\), and define its maximal absolute risk, its two output laws, and their error events by
\[
\begin{aligned}
\mathcal R(M)
&=
\sup_{P\in\mathcal P}
E_{(P^{\mathrm{obs}})^{\otimes n},M}
|M(D)-\theta(P)|,\\
\mathsf L_\eta&=MQ^{(\eta)},\\
\mathcal E_\eta
&=
\{u\in\mathbb R:\Delta/2\le |u-\eta\Delta|\},
\qquad \eta\in\{0,1\}.
\end{aligned}
\]
Each \(\mathsf L_\eta\) is a probability measure, since its dataset prior is a normalized nonempty finite mixture and \(M\) is a Markov kernel.

The target distance satisfies
\[
2(\Delta/2)=|0-\Delta|.
\]
Consequently, the triangle inequality gives
\(\mathcal E_0^c\subseteq\mathcal E_1\). The two-error testing inequality therefore yields
\[
\begin{aligned}
\mathsf L_0(\mathcal E_0)+\mathsf L_1(\mathcal E_1)
&\ge
\mathsf L_0(\mathcal E_0)+\mathsf L_1(\mathcal E_0^c)\\
&=
1+\mathsf L_0(\mathcal E_0)-\mathsf L_1(\mathcal E_0)\\
&\ge
1-d_{\mathrm{TV}}(\mathsf L_0,\mathsf L_1)
\ge\frac34.
\end{aligned}
\]
Thus
\[
\mathsf L_0(\mathcal E_0)\ge\frac38
\quad\text{or}\quad
\mathsf L_1(\mathcal E_1)\ge\frac38.
\]

In the first case take \(\eta=0\), and in the second take \(\eta=1\). The absolute loss is at least \(\Delta/2\) on the corresponding error event, so
\[
\int_{\mathbb R}|u-\eta\Delta|\,d\mathsf L_\eta(u)
\ge
\frac{\Delta}{2}\mathsf L_\eta(\mathcal E_\eta)
\ge
\frac{\Delta}{2}\frac38
=
\frac{3\Delta}{16}.
\]
% lean: CausalSmith.Stat.PrivateCateRoughdesign.absolute_loss_of_error_probability

Finite-prior averaging and the common target within each family give
\[
\begin{aligned}
\int_{\mathbb R}|u-\eta\Delta|\,d\mathsf L_\eta(u)
&=
\frac{1}{N_{\mathrm{prior},\eta}}
\sum_{\iota=1}^{N_{\mathrm{prior},\eta}}
E_{(P_{\eta,\iota}^{\mathrm{obs}})^{\otimes n},M}
|M(D)-\theta(P_{\eta,\iota})|\\
&\le
\frac{1}{N_{\mathrm{prior},\eta}}
\sum_{\iota=1}^{N_{\mathrm{prior},\eta}}\mathcal R(M)
=
\mathcal R(M).
\end{aligned}
\]
These identities and inequalities concern nonnegative extended integrals and remain valid when a risk is infinite. In either case,
\[
\frac{3\Delta}{16}\le\mathcal R(M).
\]
% lean: CausalSmith.Stat.PrivateCateRoughdesign.finite_prior_loss_le_worstRisk
% lean: CausalSmith.Stat.PrivateCateRoughdesign.finite_priors_scalar_lower

\item The numerical calibration is
\[
2^{-20}\le\frac{3}{16}\,2^{-14}.
\]
Multiplying by the nonnegative benchmark and then using the separation bound gives
\[
2^{-20}r(n,\epsilon)
\le
\frac{3}{16}\bigl(2^{-14}r(n,\epsilon)\bigr)
\le
\frac{3\Delta}{16}
\le
\mathcal R(M).
\]
Since this holds for every admissible scalar kernel, \cref{obj:def:risk-criterion} implies
\[
2^{-20}r(n,\epsilon)
\le
\inf_{M\in\mathfrak M_{n,\epsilon}(\mathbb R)}\mathcal R(M)
=
R_{n,\epsilon}.
\]
% lean: CausalSmith.Stat.PrivateCateRoughdesign.matched_risk_frontier

\item For the upper bound, use the replacement-copy Efron--Stein inequality from \citep[Theorem 3.1, p.~54]{BoucheronLugosiMassart2013} as the variance input to \cref{obj:thm:uniform-private-upper}. Apply that result with the admissible auxiliary parameters \(k=2\) and \(h=1/4\). It supplies, for the publicly tuned scalar estimator in \cref{obj:def:public-tuning},
\[
\widehat T^*\in\mathfrak M_{n,\epsilon}(\mathbb R),
\qquad
\mathcal R(\widehat T^*)\le2^{18}r(n,\epsilon).
\]
In particular, this estimator is an everywhere-defined Markov kernel satisfying pure replacement \(\epsilon\)-differential privacy. Its risk bound takes the supremum over the full class in \cref{obj:def:complete-model}, including every admissible measurable covariate density.

Using this kernel as a candidate in the minimax infimum gives
\[
R_{n,\epsilon}
=
\inf_{M\in\mathfrak M_{n,\epsilon}(\mathbb R)}\mathcal R(M)
\le
\mathcal R(\widehat T^*)
\le
2^{18}r(n,\epsilon).
\]
Together with the lower bound, this proves all the assertions.
% lean: CausalSmith.Stat.PrivateCateRoughdesign.uniform_private_upper
% lean: CausalSmith.Stat.PrivateCateRoughdesign.matched_risk_frontier
\end{enumerate}
\end{proof}
